\section{Machine-checked: \textsc{ToMatchertext}}
\label{sec:appendix:tomatchertext}

The declaration heads below summarize the Lean~4 development in
\texttt{lean/to\_matchertext.lean} that proves the write-path obligation of
\cref{sec:threat} (TCB item~iii): the total encoder's output is \emph{always}
matchertext. It reuses the base definitions of \cref{sec:appendix:lean}
(\texttt{Opener}, \texttt{Closer}, \texttt{Nonmatcher}, \texttt{MT}, and the
\texttt{depth}/\texttt{embed\_boundary} results) and the closure lemmas
\texttt{mt\_append}/\texttt{mt\_wrap} imported from the public artifact
of~\cite{matchertext}.
Definition bodies are shown in full; tactic proof bodies are elided and given in
the source file.
The encoder escapes the unmatched matchers and, by two decisions, may escape more:
an opener past a depth limit, and a matched pair. Any choice of the two keeps the
output in $L$, and \texttt{toMatchertextC} fixes them as the C encoder does: the
host's depth limit, and the pair after an introducer and tag, so that input
already shaped like an escape decodes to itself. The model is of the algorithm,
not of the C code, and does not cover \textsc{Verify}.

\begin{lstlisting}
open Classical

variable {α : Type*}         -- the alphabet Σ, kept abstract
variable (Pi : Set (α × α))  -- the set of pairs Π

-- A pending scan context: an opener awaiting its closer, with the output at its level.
abbrev Ctx (α : Type*) := α × List α

-- End of input: every opener still pending is unmatched, so it is escaped.
noncomputable def unwind (esc : α → List α) : List (Ctx α) → List α → List α
  | [], out => out
  | (o, b) :: st, out => unwind esc st (b ++ esc o ++ out)

-- One left-to-right pass. `st` holds the pending openers, `out` the output at the current
-- level. A closer is matched only if it pairs with the top; else it is unmatched and escaped.
-- Two decisions may escape more: `cap st` escapes an opener instead of nesting it (a depth
-- limit), and `split b out` escapes both members of a matched pair, given the output `b`
-- before its opener and its body `out`. Any choice of either keeps the output matchertext.
noncomputable def encAux (esc : α → List α) (cap : List (Ctx α) → Bool)
    (split : List α → List α → Bool) :
    List (Ctx α) → List α → List α → List α
  | st, out, [] => unwind esc st out
  | st, out, x :: xs =>
      if Opener Pi x then
        if cap st then encAux esc cap split st (out ++ esc x) xs
        else encAux esc cap split ((x, out) :: st) [] xs
      else if Closer Pi x then
        match st with
        | (o, b) :: st' =>
            if (o, x) ∈ Pi then
              if split b out then encAux esc cap split st' (b ++ esc o ++ out ++ esc x) xs
              else encAux esc cap split st' (b ++ [o] ++ out ++ [x]) xs
            else encAux esc cap split st (out ++ esc x) xs
        | [] => encAux esc cap split st (out ++ esc x) xs
      else encAux esc cap split st (out ++ [x]) xs

-- The minimal encoder: only unmatched matchers are escaped.
noncomputable def toMatchertext (esc : α → List α) (v : List α) : List α :=
  encAux Pi esc (fun _ => false) (fun _ _ => false) [] [] v

-- The C encoder (sqlite/src/matchertext.c): nesting stops one level below the host's
-- limit `n`, and a matched pair right after an introducer `bs` and a tag `t` with
-- `tag t` is escaped, so that input already shaped like an escape decodes to itself.
noncomputable def toMatchertextC (esc : α → List α) (n : Nat) (bs : α) (tag : α → Prop)
    (v : List α) : List α :=
  encAux Pi esc (fun st => decide (n ≤ st.length + 1))
    (fun b out => decide (out = [] ∧ ∃ p t, b = p ++ [bs, t] ∧ tag t)) [] [] v

-- Lemma: the escapes this implementation ships, a backslash and a tag followed
-- by a matched pair as in \o(), are matchertext by nesting rather than flatness.
theorem esc_mt_paired
    (hdisj : ∀ x y z, (x, y) ∈ Pi → (z, x) ∈ Pi → False)
    {bs t o c : α}
    (hbs : Nonmatcher Pi bs) (ht : Nonmatcher Pi t) (hp : (o, c) ∈ Pi) :
    MT Pi [bs, t, o, c] ∧ depth Pi [bs, t, o, c] = 0 := by ...

-- Invariant (the crux): every accumulated segment stays matchertext, whatever the
-- decisions `cap` and `split` choose.
theorem encAux_mt (esc : α → List α)
    (hesc : ∀ c, MT Pi (esc c) ∧ depth Pi (esc c) = 0)
    (cap : List (Ctx α) → Bool) (split : List α → List α → Bool)
    (st : List (Ctx α)) (hst : ∀ p ∈ st, MT Pi p.2)
    (out : List α) (hout : MT Pi out) (v : List α) :
    MT Pi (encAux Pi esc cap split st out v) := by ...

-- Theorem (output ∈ L): any string at all, once encoded, is matchertext.
theorem toMatchertext_mt (esc : α → List α)
    (hesc : ∀ c, MT Pi (esc c) ∧ depth Pi (esc c) = 0) (v : List α) :
    MT Pi (toMatchertext Pi esc v) := ...

-- Theorem: the same for the C encoder's decisions.
theorem toMatchertextC_mt (esc : α → List α)
    (hesc : ∀ c, MT Pi (esc c) ∧ depth Pi (esc c) = 0)
    (n : Nat) (bs : α) (tag : α → Prop) (v : List α) :
    MT Pi (toMatchertextC Pi esc n bs tag v) := ...

-- Corollary (write-path payoff): reading to matcher balance recovers exactly the encoded
-- value, for a freely chosen v. embed_boundary assumes MT Pi m, which toMatchertext_mt supplies.
theorem embed_boundary_encoded
    (hdisj : ∀ x y z, (x, y) ∈ Pi → (z, x) ∈ Pi → False)
    (esc : α → List α)
    (hesc : ∀ c, MT Pi (esc c) ∧ depth Pi (esc c) = 0)
    (v : List α) {c : α} (hc : Closer Pi c) :
    (∀ p, p <+: toMatchertext Pi esc v → 0 ≤ depth Pi p)
      ∧ depth Pi (toMatchertext Pi esc v) = 0
      ∧ depth Pi (toMatchertext Pi esc v ++ [c]) = -1 := ...
\end{lstlisting}
